\chapter{Variable Definitions and Supplementary Notes}
\label{chap:appa-variable-notes-en}

\section{Definitions of Core Variables}
\label{sec:appa-variable-definitions-en}

This appendix summarizes the key variables used in the empirical chapters and provides a compact example of how appendix content can support the main text. Supplementary robustness results are reported in Appendix~\ref{chap:appb-robustness-tables-en}. Table~\ref{tab:appa-core-variables-en} lists the main variables, their operationalization, and their analytical purpose, while equation~\eqref{eq:appa-monitoring-index-en} illustrates a simple monitoring intensity index constructed from standardized components.

\begin{table}[htbp]
    \centering
    \caption{Illustrative Definitions of Core Variables}
    \label{tab:appa-core-variables-en}
    \begin{tabular}{p{3.3cm}p{4.1cm}p{5.5cm}}
        \toprule
        Variable & Measurement & Description \\
        \midrule
        shirking\_score & Survey index from 1 to 5 & Higher values indicate more frequent non-task activities during working hours. \\
        monitor\_intensity & Standardized monitoring count & Combines spot checks, system traces, and supervisor reviews at the monthly level. \\
        incentive\_pay & Share of performance pay & Measures the proportion of variable compensation in total income. \\
        task\_complexity & Task complexity score & Captures task diversity, knowledge intensity, and cross-team coordination needs. \\
        \bottomrule
    \end{tabular}
\end{table}

For smoke-test purposes, the appendix also contains a numbered equation:
\begin{equation}
    \label{eq:appa-monitoring-index-en}
    M_i = \frac{1}{3}\left(z(\text{spotcheck}_i) + z(\text{logtrace}_i) + z(\text{managerreview}_i)\right).
\end{equation}

The main text may therefore refer readers to Table~\ref{tab:appa-core-variables-en} and equation~\eqref{eq:appa-monitoring-index-en} when describing the construction of the monitoring variable.
